% ============================================================================
%  section_v_sensitivity.tex
%  NEW Section V-F — paste immediately BEFORE \section{Conclusion}
% ============================================================================

\subsection{Parameter Sensitivity Analysis}
\label{sec:sensitivity}

We introduced four tuning parameters---$\alpha$, $\beta$, $\lambda$, and
$\kappa$---with distinct roles in the confidence filter and risk modulation
(Sections~III--IV).
Fig.~\ref{fig:sensitivity} reports a one-at-a-time parameter sweep on
Scenario~A with the adaptive controller; all other parameters remain at their
nominal values used throughout Section~V.

\begin{figure}[t]
\centering
\includegraphics[width=\linewidth]{fig_sensitivity.pdf}
\caption{Parameter sensitivity on Scenario~A (adaptive controller).
Each panel varies one parameter while holding the others at nominal values
($\alpha=0.30$, $\beta=0.33$, $\lambda=3$, $\kappa=0.6$; dashed vertical
lines). Mean back-off varies smoothly; the nominal setting lies in the
interior of a flat region for each parameter, indicating robustness to
moderate retuning.}
\label{fig:sensitivity}
\end{figure}

The evidence gain~$\alpha$ and prior-anchor gain~$\beta$ control the
erosion--recovery asymmetry of the MAP filter (Proposition~1): larger
$\alpha$ reacts faster to prediction errors; larger~$\beta$ damps transient
confidence swings.
Risk-modulation parameters~$\lambda$ and~$\kappa$ affect back-off magnitude
monotonically, consistent with Theorem~\ref{thm:feasible}.
Across the tested ranges, no setting produced a safety violation on the
representative Scenario~A trajectory; the severe-regime Monte Carlo study
(Table~\ref{tab:montecarlo}) remains the distribution-level stress test for
calibration.
